\Needspace{18\baselineskip}
\subsection{Models over a curve}
\label{subsec:curve-fibration-polar-models}

The polar metric on a curve also describes models on a compact
manifold that fibers smoothly over it. A background pulled back
from the curve has no vertical curvature, so every potential
descends. Gauduchon fiber volumes are constant; consequently the
normalized polar trace recovers the curve metric without choosing
a product structure.

Let $p:X\to C$ be a surjective holomorphic submersion with connected
fibers, where $X$ is a connected compact complex manifold of
dimension $n\geq2$ and $C$ is a connected compact Riemann surface.
Fix a smooth positive $(1,1)$-form $\beta$ on $C$, put
$A=\int_C\beta$ and $\theta=p^*\beta$. For
$v\in\PSH(C,\beta)$ write
\[
 T_v=\beta+\ddc v,\qquad
 \mu_v=\mathbf1_{\{v=-\infty\}}T_v.
\]
Thus $\mu_v$ is the polar curvature of the curve potential.

\Needspace{14\baselineskip}
\begin{lemma}\label{fibration:curve-envelope-pullback}\index{Envelope!curve submersion}
For $v\in\PSH(C,\beta)$ and $U=v\circ p$,
\[
 \begin{aligned}
 P_\theta[U]&=P_\beta[v]\circ p,\\
 P_\theta[U]_{\mathcal I}&=P_\beta[v]_{\mathcal I}\circ p.
 \end{aligned}
\]
Pullback is a bijection between the ordinary model spaces and
between their ideal subspaces. Every $U\in\PSH(X,\theta)$ has zero
extended Monge--Amp\`ere measure. In particular all ordinary types
in this background have zero upper and lower singularity volumes.
\end{lemma}
\begin{proof}
By \cref{prop:psh-descent-submersion}, every potential in the
pulled-back background is uniquely $v\circ p$. Surjectivity of $p$
preserves the supremum and the ordinary singularity order.

By \cref{lem:submersion-multiplier-pullback}, multiplier ideals
commute with this pullback. In local product coordinates, a curve
ideal $(z^r)$ extends to the product ideal generated by $z^r$;
equality of these extended ideals therefore implies equality on
the base. Thus ordinary and ideal competitor
families correspond bijectively. Upper regularization commutes
with pullback in these coordinates, giving both envelope identities
and the assertions about fixed points.

For the mass assertion, $\max(U,-j)=\max(v,-j)\circ p$.
Locally write $\beta=\ddc a$ with $a$ smooth. The bounded
plurisubharmonic function $(a+\max(v,-j))\circ p$ has smooth
decreasing local approximations depending only on $z$.
Their complex Hessians have rank at most one, so their $n$th
powers vanish. Bounded product continuity gives
$(\theta+\ddc\max(U,-j))^n=0$. Each non-pluripolar truncation is
therefore zero, proving the measure assertion. Since every
potential in the background descends, the same conclusion holds
for every representative of every ordinary type.
\end{proof}

Fix a Gauduchon form $\gamma$ on $X$. For a curve measure $\mu$,
the positive current $p^*\mu$ is defined in submersion coordinates
by pulling back its $(1,1)$-current; a signed measure is treated
by linearity. Write
\[
 L_\gamma(\mu)=p^*\mu\wedge\gamma^{n-1}.
\]
This is a finite signed measure on $X$ when $\mu$ is finite.

\Needspace{14\baselineskip}
\begin{lemma}\label{fibration:gauduchon-fiber-volume}\index{Gauduchon metric!constant fiber volume}
The fiber integral
\[
 \kappa_\gamma(t)=\int_{p^{-1}(t)}\gamma^{n-1}
\]
is a positive constant. For every finite signed measure $\mu$ on $C$,
\[
 p_*L_\gamma(\mu)=\kappa_\gamma\mu,\qquad
 \|L_\gamma(\mu)\|_{\mathrm{TV}}
             =\kappa_\gamma\|\mu\|_{\mathrm{TV}}.
\]
\end{lemma}
\begin{proof}
Smooth local trivializations of the proper submersion show that
the fiber integral is smooth; it is positive since $\gamma$ is
positive and every fiber has dimension $n-1$. For a smooth test
function $h$ on $C$, fiber integration and Stokes' theorem give
\[
 \int_C\kappa_\gamma\ddc h
 =\int_X\gamma^{n-1}\wedge p^*\ddc h
 =\int_X(h\circ p)\ddc(\gamma^{n-1})=0.
\]
Thus $\kappa_\gamma$ is harmonic and hence constant on the compact
connected curve. Fubini's theorem in product charts proves
$p_*L_\gamma(\mu)=\kappa_\gamma\mu$ first for positive measures
and then for signed measures.

For the variation identity, let $\mu=\mu^+-\mu^-$ be the Jordan
decomposition. Choose disjoint Borel carriers for $\mu^+$ and
$\mu^-$. Their inverse images carry the two positive measures
$L_\gamma(\mu^+)$ and $L_\gamma(\mu^-)$ and are disjoint, so these
are the Jordan parts of $L_\gamma(\mu)$. Their masses are
$\kappa_\gamma\mu^+(C)$ and $\kappa_\gamma\mu^-(C)$ by the
pushforward identity, which gives the norm equality.
\end{proof}

Let $\mathcal M(\theta)$ be the ordinary model space. For
$U=v\circ p\in\mathcal M(\theta)$, its polar trace measure is
\[
 \lambda_U^\gamma
 =\mathbf1_{\{U=-\infty\}}
       (\theta+\ddc U)\wedge\gamma^{n-1}
 =L_\gamma(\mu_v).
\]
The last identity follows in local product coordinates, including
when the polar measure has a nonatomic part. Define
\[
 \rho_\theta(U,V)=\kappa_\gamma^{-1}
                  \|\lambda_U^\gamma-\lambda_V^\gamma\|_{\mathrm{TV}}.
\]

\Needspace{14\baselineskip}
\begin{proposition}\label{fibration:polar-model-metric}\index{Metric completeness!curve submersion}
The metric $\rho_\theta$ is independent of the Gauduchon form
$\gamma$. Pullback is an isometry from the complete geodesic
curve model space onto $(\mathcal M(\theta),\rho_\theta)$, which
has diameter $2A$. The ideal models form a closed geodesically
convex subspace equal to the metric closure of the models with
analytic singularities. Ideal projection is its unique
nearest-point retraction and is $1$-Lipschitz. If $U=v\circ p$ is
a model, its distance to that subspace is the nonatomic polar
mass of $v$.

Metric convergence implies $L^1(X)$ convergence for every smooth
positive volume form. The upper-volume comparison on
$\mathcal M(\theta)$ vanishes identically.
\end{proposition}
\begin{proof}
The preceding two lemmas give
\[
 \rho_\theta(v\circ p,w\circ p)
       =\|\mu_v-\mu_w\|_{\mathrm{TV}}
       =d_\beta(v,w).
\]
This proves independence of $\gamma$ and the isometry.
Completeness and geodesics follow from
\cref{cor:curve-model-completeness}. Every polar measure has mass
at most $A$, so the diameter is at most $2A$; the models of
$A\delta_s$ and $A\delta_t$ for distinct $s,t\in C$ attain it.
The envelope identity identifies ideal projection, so the closed
ideal subspace, nearest-point property, Lipschitz bound and
distance formula follow from \cref{prop:curve-analytic-distance-defect}.

Analytic singularities pull back in product coordinates.
Conversely a pulled-back model with analytic singularities
restricts to a local holomorphic section through each point of
$C$, and the restriction has analytic singularities by
\cref{prop:analytic-type-calculus}. It is not identically
$-\infty$, since the descended potential is a curve potential.
Thus the analytic models correspond, and their metric closure
is the ideal subspace: every analytic model is ideal by
\cref{prop:analytic-envelope-type}, while the analytic curve
models of positive curve mass are dense in that subspace by
\cref{cor:curve-model-completeness}. Their pullbacks have zero
top-degree mass, as do all models in this background.

For a smooth positive volume form $\mathrm dV$ on $X$, its
pushforward is a smooth positive volume form on $C$, bounded
above by $C_0\beta$ for some $C_0>0$. Consequently
\[
 \int_X|(v_j-v)\circ p|\,\mathrm dV
      \leq C_0\int_C|v_j-v|\,\beta.
\]
Apply \cref{thm:curve-model-capacity-convergence} on the curve.
Finally every potential in the background has zero extended
mass, including the pairwise maximum, by the first lemma.
The defining upper-volume comparison therefore vanishes.
\end{proof}

\Needspace{14\baselineskip}
\begin{corollary}\label{fibration:all-obstacle-criterion}\index{Envelope!all-obstacle curve submersion criterion}
Let $U=v\circ p$ and $R=P_\theta[U]_{\mathcal I}$. The following
are equivalent: $P_\theta[U]=P_\theta[U]_{\mathcal I}$;
$\mu_v$ is atomic; $U\in\mathcal E_{\mathrm{env}}(X,\theta,R)$;
and the ordinary and ideal envelopes of $U$ agree below every
smooth real obstacle on $X$.
\end{corollary}
\begin{proof}
Put $r=P_\beta[v]_{\mathcal I}$, so $R=r\circ p$ by the first
lemma. The same lemma makes the first equality equivalent to
$P_\beta[v]=r$. By \cref{cor:curve-arbitrary-background-criterion},
this holds exactly when $\mu_v$ is its atomic part, and exactly
when $v\in\mathcal E_{\mathrm{env}}(C,\beta,r)$.

Surjectivity and descent preserve the ordinary order and give
a bijection between all stretching competitors: for every
$b\geq1$, the inequality
\[
 b^{-1}h+(1-b^{-1})r\leq v
\]
holds on $C$ if and only if its pullback holds on $X$.
This remains true at poles because only convex combinations
are used. Thus the two relative-class conditions are equivalent.
Membership upstairs and \cref{prop:ideal-obstacle-reprojection}
give equality below every smooth obstacle. Conversely that
equality includes the zero obstacle and implies the first
condition, completing the equivalence.
\end{proof}

The top-degree measure is zero for every potential, yet the
first polar trace separates every pair of models. Here the
one-dimensional quotient supplies all the singularity data.

\Needspace{18\baselineskip}
\subsection{Hopf surface fibers and unbounded ambient mass}
\label{subsec:hopf-surface-fibration}

The complete model space of a curve submersion can occur on a
manifold whose ambient Hermitian masses are unbounded. Blowing
up the elliptic base locus of the projection $[z_0:z_1]$ on the
scalar Hopf threefold turns this projection into a submersion
with Hopf surface fibers. The rank-one background then retains
only the curve singularities,
while the smooth potentials of the Hopf threefold still force
unbounded mass in a positive ambient background.

\Needspace{16\baselineskip}
\begin{example}\label{fibration:unbounded-hopf-threefold}\index{Hopf manifold!surface fibration}\index{Metric completeness!Hopf surface fibration}\index{Bounded mass property!curve submersion counterexample}
There is a connected compact complex threefold $Z$ with a
surjective holomorphic submersion $p:Z\to\mathbb P^1$ whose
fibers are scalar Hopf surfaces. The manifold $Z$ does not have
the bounded mass property. For the pullback
$\theta=p^*\omega_{\FS}$ of the Fubini--Study form of area one,
the entire ordinary model space has the complete geodesic polar
metric of \cref{fibration:polar-model-metric}, with diameter two.
Its ideal subspace is the analytic closure, with the unique
nearest-point, $1$-Lipschitz ideal retraction. All potentials
in this background have zero extended Monge--Amp\`ere measure,
yet it contains analytic ideal models at polar distance one.
\end{example}
\begin{proof}
Put $\lambda=\mathrm e^{-1}$ and consider the vector bundle
\[
 \mathcal V=\mathcal O_{\mathbb P^1}(-1)
                  \oplus\mathcal O_{\mathbb P^1}.
\]
Let $\mathcal V^\times$ be its total space with the zero section
removed, and define
\[
 Z=\mathcal V^\times/\langle v\mapsto\lambda v\rangle.
\]
The bundle projection descends to a holomorphic map
$p:Z\to\mathbb P^1$.
A Hermitian norm on the bundle satisfies
$\|\lambda^kv\|=\lambda^k\|v\|$. The action is free and proper:
on compact subsets of $\mathcal V^\times$ the norms have positive
lower and finite upper bounds, so only finitely many powers can
move one such subset to meet another. Thus the quotient is a
complex threefold. The annulus bundle
$\{\lambda\leq\|v\|\leq1\}$ is compact over $\mathbb P^1$ and
meets every orbit, so $Z$ is compact. Bundle charts identify
$p^{-1}(U)$ with
\[
 U\times\bigl((\mathbb C^2\setminus\{0\})/
                    \langle w\mapsto\lambda w\rangle\bigr).
\]
Hence $p$ is a holomorphic submersion with connected compact
Hopf surface fibers, and $Z$ is connected.

The tautological inclusion into $\mathbb C^2$ identifies
$\mathcal V^\times$ with the incidence manifold
\[
 \left\{([s:t],z)\in\mathbb P^1\times
      (\mathbb C^3\setminus\{0\}):\ t z_0=s z_1\right\}.
\]
The last coordinate $z_2$ comes from the trivial summand. This
is the blowup of $\mathbb C^3\setminus\{0\}$ along
$L=\{z_0=z_1=0,\ z_2\ne0\}$: in the chart $s\ne0$, write $a=t/s$ and
$z_1=az_0$, while in the chart $t\ne0$ write $b=s/t$ and
$z_0=bz_1$. These are the two usual blowup charts of the ideal
$(z_0,z_1)$. The scalar contraction preserves the incidence
relation and commutes with projection to $z$. It therefore
descends to a modification
\[
 \mu:Z\longrightarrow
 H=(\mathbb C^3\setminus\{0\})/
                \langle z\mapsto\lambda z\rangle.
\]
The map is an isomorphism away from the elliptic curve
$L/\langle\lambda\rangle$. Local quotient charts show that it
is its blowup; in particular it has degree one.

Take the Gauduchon form $\omega$ and smooth potentials $V_j$ on
$H$ from \cref{thm:hopf-unbounded}. Since $Z$ is compact, a
sufficiently large multiple $\Omega$ of any Hermitian form satisfies
$\Omega-\mu^*\omega>0$. Put $W_j=V_j\circ\mu$. Then
\[
 \Omega+\ddc W_j
 =\Omega-\mu^*\omega+\mu^*(\omega+\ddc V_j)>0.
\]
Positivity of smooth mixed products and the degree-one change
of variables off the exceptional divisor give
\[
 \begin{aligned}
 \int_Z(\Omega+\ddc W_j)^3
 &\geq\int_Z\mu^*(\omega+\ddc V_j)^3\\
 &=\int_H(\omega+\ddc V_j)^3\longrightarrow+\infty.
 \end{aligned}
\]
The exceptional divisor and its image have zero measure for
smooth top-degree forms, so they do not affect the equality.
Every $W_j$ is smooth and bounded, so $v_+(\Omega)=+\infty$.
By \cref{prop:bmp-metric}, $Z$ does not have the bounded mass
property.

Apply \cref{fibration:curve-envelope-pullback,fibration:polar-model-metric} to $p$ and the area-one form
$\omega_{\FS}$. This gives all the assertions about the full
model space. To see the stated distance explicitly, fix
$x\in\mathbb P^1$ and let $g_x$ be its normalized Green model,
with $\omega_{\FS}+\ddc g_x=\delta_x$. The models $0$ and
$g_x\circ p$ are ordinary and ideal, and the latter has a
logarithmic pole of coefficient one along the smooth reduced
fiber $p^{-1}(x)$. Thus both have analytic singularities. Their
base polar measures are $0$ and $\delta_x$, so their polar
distance is one, although their top-degree measures vanish.
\end{proof}
